% Preamble for transcriptions of documents of the Bourbaki archive
% (Archives Henri Poincaré, archives-bourbaki.ahp-numerique.fr), issue #50.
%
% The fonds' preamble, whole: the same apparatus macros, tikz-cd set-up and
% watermark, so that one renderer can read both one day. Only the head changes:
% it names the rédaction and the archive whose facsimile it was read from.
%
% A transcription sits one directory below transcripts-bourbaki/, as
% transcripts-bourbaki/<number>/batch-NN.fr.tex, and opens with
%   % Preamble for transcriptions of documents of the Bourbaki archive
% (Archives Henri Poincaré, archives-bourbaki.ahp-numerique.fr), issue #50.
%
% The fonds' preamble, whole: the same apparatus macros, tikz-cd set-up and
% watermark, so that one renderer can read both one day. Only the head changes:
% it names the rédaction and the archive whose facsimile it was read from.
%
% A transcription sits one directory below transcripts-bourbaki/, as
% transcripts-bourbaki/<number>/batch-NN.fr.tex, and opens with
%   % Preamble for transcriptions of documents of the Bourbaki archive
% (Archives Henri Poincaré, archives-bourbaki.ahp-numerique.fr), issue #50.
%
% The fonds' preamble, whole: the same apparatus macros, tikz-cd set-up and
% watermark, so that one renderer can read both one day. Only the head changes:
% it names the rédaction and the archive whose facsimile it was read from.
%
% A transcription sits one directory below transcripts-bourbaki/, as
% transcripts-bourbaki/<number>/batch-NN.fr.tex, and opens with
%   % Preamble for transcriptions of documents of the Bourbaki archive
% (Archives Henri Poincaré, archives-bourbaki.ahp-numerique.fr), issue #50.
%
% The fonds' preamble, whole: the same apparatus macros, tikz-cd set-up and
% watermark, so that one renderer can read both one day. Only the head changes:
% it names the rédaction and the archive whose facsimile it was read from.
%
% A transcription sits one directory below transcripts-bourbaki/, as
% transcripts-bourbaki/<number>/batch-NN.fr.tex, and opens with
%   \input{../preamble/bourbaki}

\input{../../transcripts/preamble/grothendieck}

\newcommand{\redaction}[1]{\folder{#1}}

\renewcommand{\grothTitleBlock}{%
  \begin{center}
    {\large\bfseries Archives Bourbaki --- rédaction n\textsuperscript{o}\,\grothFolder}\\[2pt]
    \ifx\grothFoldertitle\empty\else{\itshape\grothFoldertitle}\\[4pt]\fi
    {\small pages \grothFirst--\grothLast
      \ifx\grothBatch\empty\else\ (lot \grothBatch)\fi}\\[2pt]
    \ifx\grothDating\empty\else
      {\small\color{gray!60!black}Datation de l'archive : \grothDating}\\[2pt]
    \fi
    \ifx\grothWatermark\empty\else
      {\small\bfseries\color{red!45!gray}\grothWatermark}\\[2pt]
    \fi
    {\footnotesize\color{gray!60!black}Transcrit du fac-similé publié par les Archives Henri Poincaré
     (archives-bourbaki.ahp-numerique.fr).\\
     Lectures douteuses soulignées, passages illisibles marqués [\,\ldots\,], ajouts entre crochets.}
  \end{center}
  \vspace{1em}}

\endinput


\input{../../transcripts/preamble/grothendieck}

\newcommand{\redaction}[1]{\folder{#1}}

\renewcommand{\grothTitleBlock}{%
  \begin{center}
    {\large\bfseries Archives Bourbaki --- rédaction n\textsuperscript{o}\,\grothFolder}\\[2pt]
    \ifx\grothFoldertitle\empty\else{\itshape\grothFoldertitle}\\[4pt]\fi
    {\small pages \grothFirst--\grothLast
      \ifx\grothBatch\empty\else\ (lot \grothBatch)\fi}\\[2pt]
    \ifx\grothDating\empty\else
      {\small\color{gray!60!black}Datation de l'archive : \grothDating}\\[2pt]
    \fi
    \ifx\grothWatermark\empty\else
      {\small\bfseries\color{red!45!gray}\grothWatermark}\\[2pt]
    \fi
    {\footnotesize\color{gray!60!black}Transcrit du fac-similé publié par les Archives Henri Poincaré
     (archives-bourbaki.ahp-numerique.fr).\\
     Lectures douteuses soulignées, passages illisibles marqués [\,\ldots\,], ajouts entre crochets.}
  \end{center}
  \vspace{1em}}

\endinput


\input{../../transcripts/preamble/grothendieck}

\newcommand{\redaction}[1]{\folder{#1}}

\renewcommand{\grothTitleBlock}{%
  \begin{center}
    {\large\bfseries Archives Bourbaki --- rédaction n\textsuperscript{o}\,\grothFolder}\\[2pt]
    \ifx\grothFoldertitle\empty\else{\itshape\grothFoldertitle}\\[4pt]\fi
    {\small pages \grothFirst--\grothLast
      \ifx\grothBatch\empty\else\ (lot \grothBatch)\fi}\\[2pt]
    \ifx\grothDating\empty\else
      {\small\color{gray!60!black}Datation de l'archive : \grothDating}\\[2pt]
    \fi
    \ifx\grothWatermark\empty\else
      {\small\bfseries\color{red!45!gray}\grothWatermark}\\[2pt]
    \fi
    {\footnotesize\color{gray!60!black}Transcrit du fac-similé publié par les Archives Henri Poincaré
     (archives-bourbaki.ahp-numerique.fr).\\
     Lectures douteuses soulignées, passages illisibles marqués [\,\ldots\,], ajouts entre crochets.}
  \end{center}
  \vspace{1em}}

\endinput


\input{../../transcripts/preamble/grothendieck}

\newcommand{\redaction}[1]{\folder{#1}}

\renewcommand{\grothTitleBlock}{%
  \begin{center}
    {\large\bfseries Archives Bourbaki --- rédaction n\textsuperscript{o}\,\grothFolder}\\[2pt]
    \ifx\grothFoldertitle\empty\else{\itshape\grothFoldertitle}\\[4pt]\fi
    {\small pages \grothFirst--\grothLast
      \ifx\grothBatch\empty\else\ (lot \grothBatch)\fi}\\[2pt]
    \ifx\grothDating\empty\else
      {\small\color{gray!60!black}Datation de l'archive : \grothDating}\\[2pt]
    \fi
    \ifx\grothWatermark\empty\else
      {\small\bfseries\color{red!45!gray}\grothWatermark}\\[2pt]
    \fi
    {\footnotesize\color{gray!60!black}Transcrit du fac-similé publié par les Archives Henri Poincaré
     (archives-bourbaki.ahp-numerique.fr).\\
     Lectures douteuses soulignées, passages illisibles marqués [\,\ldots\,], ajouts entre crochets.}
  \end{center}
  \vspace{1em}}

\endinput
